\textbf{结论名称：}圆内接四边形的判定（四点共圆的等价条件）

\textbf{适用条件：}平面内不共线的四点 \(A,B,C,D\)，或按序连成的四边形 \(ABCD\)。

\textbf{变量说明：}四边形 \(ABCD\) 的内角记作 \(\angle A,\angle B,\angle C,\angle D\)；若延长 \(BC\) 至 \(E\)，则 \(\angle DCE\) 为外角。

\textbf{核心公式：}
\[
\boxed{A,B,C,D\text{ 四点共圆} \iff \angle A+\angle C=180^\circ \iff \angle B+\angle D=180^\circ}
\]
等价形式：\(\angle DCE=\angle A\)（外角等于内对角）；同侧等角 \(\angle ADB=\angle ACB\)（对同底 \(AB\)）。

\textbf{使用场景：}几何证明中判定四点共圆，进而运用圆周角定理、相交弦定理；解析几何中将点坐标代入统一的圆的方程验证共圆。

\textbf{简单例子：}在四边形 \(ABCD\) 中，若 \(\angle A=80^\circ,\;\angle C=100^\circ\)，则 \(\angle A+\angle C=180^\circ\)，故 \(A,B,C,D\) 四点共圆，可在该外接圆上继续研究其他角的关系。